\chapter{李群}

本章中，K 表示以下三者之一：实数的带赋值域 R、复数的带赋值域 C，或非离散完备超度量交换域。我们假定从 § 4 起 K 的特征为 0，§ 6 中 K = R 或 C，§ 7 中 K 为超度量域。除非另有说明，所考虑的所有流形、代数和向量空间都定义在 K 上。回忆一下，当我们说一个 \(C^{r}\) 类流形时，r ∈ N\_K，即若 K ≠ R 则 r = ω，而若 K = R 则 l ≤ r ≤ ω。

关于范数、可赋范空间和赋范空间的约定采用《微分与解析流形》R 中的约定。

回忆一下，K 上的可赋范代数是 K 上的一个（不一定结合的）代数 A，带有满足下列性质的拓扑 T：

\begin{enumerate}
\def\labelenumi{(\arabic{enumi})}
\item
  \(\mathcal{T}\) 可由范数定义；
\item
  映射 \((x,y)\mapsto xy\) 从 \(\mathbf{A}\times \mathbf{A}\) 到 \(\mathbf{A}\) 是连续的。
\end{enumerate}

A 的双连续自同构群记为 \(\mathrm{Aut(A)}\)。每个有限维 \(K\) 上的代数都是带有典范拓扑的可赋范代数。\(K\) 上的赋范代数是带有范数的 \(K\) 上代数 A，满足范数不等式 \(\| xy\| \leqslant \| x\|\| y\|\) 对 A 中所有 \(x,y\) 都成立；赋予由该范数定义的拓扑后，代数 A 是可赋范代数。若 A 是可赋范代数，则 A 上存在一个定义其拓扑并使其成为赋范代数的范数。

若 G 是群，则 \(e_{G}\)，或简写为 e，表示 G 的单位元。对 \(g \in G\)，\(\gamma(g)\)、\(\tilde{\delta}(g)\) 和 \(\operatorname{Int}(g)\) 分别表示 G 到 G 的映射 \(g' \mapsto gg'\)、\(g' \mapsto g'g^{-1}\) 和 \(g' \mapsto gg'g^{-1}\)。若 f 是从 G 到集合 E 的映射，则 f 表示从 G 到 E 的映射 \(g \mapsto f(g^{-1})\)。

